\section{Experimental methodology}\label{sec:setup}

The evaluation addresses seven questions in four groups: whether streamed, parallel and native output equal the reference (RQ1), how the system scales (RQ2), how well it detects on held-out data (RQ3--RQ5) and how well it detects on real data (RQ6--RQ7). \textbf{RQ1}: do the implementation's streamed, parallel and native outputs match the reference in practice, and how often does chunking without the guarantee leak? \textbf{RQ2}: how do memory and throughput scale? \textbf{RQ3}: how well do the pattern tiers detect credentials compared with purpose-built scanners? \textbf{RQ4}: how well does SafeStream detect personal data compared with Presidio, with and without a named-entity model? \textbf{RQ5}: how much protection does it give at the level of a whole document, and how much of that depends on rules written after seeing the benchmark? \textbf{RQ6}: how well do the systems detect personal data in real system logs? \textbf{RQ7}: how well do they detect personal data in real free text annotated by people other than the authors?

\subsection{Corpora and splits}
Table~\ref{tab:corpora} lists the test splits. Each corpus was split once into development and test halves with a fixed seed. The manifests list document identifiers and the SHA-256 of the source data, and the loader refuses a source that no longer matches. Development, error analysis and tuning used development halves only, with the exception of RedactionBench described below; test halves were scored after the code was frozen, with one disclosed exception (Section~\ref{sec:fp}).
\begin{itemize}
\item \textbf{presidio-research} \cite{presidioresearch} (synthetic): the dataset distributed with Presidio's evaluation toolkit (revision 0cb3650), with person, location, organisation, telephone, card, date and other labels, generated from templates. The 750 test documents come from 200 templates (median 3.5 documents per template).
\item \textbf{secrets} (semi-synthetic): the labelled examples in the Nosey Parker rule set \cite{noseyparker} (revision 2e6e7f3) as positives, with gold spans given by each rule's own capture groups, plus 200 generated non-secrets (UUIDs, commit hashes, SHA-256 digests, base64 image payloads, minified JavaScript) as gold-free documents. The split is by rule, so the test half contains vendor formats absent from the development half; its 216 positives come from 96 rules. The positives are short examples written by rule authors, not credentials observed in real artefacts.
\item \textbf{RedactionBench} \cite{brynjolfsson2026redactionbench}: 200 documents across eleven domains, a majority seeded from real-world sources according to its authors, with character spans labelled \emph{mandatory} or \emph{contextual}. The test half holds between 5 and 17 documents per domain. All 200 documents were used for error analysis and rule writing before the split existed, so its test half is not held out (Section~\ref{sec:contam}).
\item \textbf{OpenSSH logs} (real): the 2,000-line OpenSSH sample of Loghub \cite{zhu2023loghub}, real \texttt{sshd} lines from a server under attack, one document per line. Gold comes from labels on the 27 log templates that Loghub's parser assigns: the value of a slot that names a remote host is \texttt{IP\_ADDRESS} when it is a valid IP address (reverse-DNS hostnames are excluded), and the value of a slot that names an account, and any account name written into a template, is \texttt{USERNAME}. Applying the labels gives 1,349 IP addresses and 1,142 user names; 148 lines contain neither. The labels were drafted from the template text with an AI assistant and checked by the first author, and the corpus was frozen as test-only (commit \texttt{099cc9c}) before any system was run on it. No SafeStream rule was written with these logs in view.
\item \textbf{Linux logs} (real, held out): Loghub's 2,000-line Linux sample (sshd authentication through PAM, \texttt{su}, \texttt{login}, \texttt{ftpd}, \texttt{klogind} and the kernel), one document per line, labelled under the same rule. That gives 1,258 IP addresses and 618 user names; 407 lines contain neither. Its labels were drafted only after the user-name tier of Section~\ref{sec:logusers} had been committed (commit \texttt{c699ec6}), checked by the first author, and frozen as test-only (commit \texttt{7f87a77}) before any system was run on it, so the corpus is held out from that tier and from every other rule.
\item \textbf{TAB} (real, human gold, held out): the test split of the Text Anonymization Benchmark \cite{pilan2022tab} (repository commit \texttt{558e09e}, MIT licence), 127 English judgments of the European Court of Human Rights annotated by one to ten trained annotators each (555 document--annotator pairs). Annotators marked every mention that identifies the applicant directly (\textsc{direct}) or in combination with other information (\textsc{quasi}), with types such as person, code, date, organisation and location; these mentions are gold. No rule of the evaluated configurations was written with TAB in view. Its manifest and its scoring protocol (metrics, systems and planned comparisons) were committed before any system was run on it (commit \texttt{17b80c4}). TAB's two other splits serve the free-text configuration of Section~\ref{sec:prose}: its development split (127 judgments) is the only TAB data used to design it, and its training split (1,014 judgments, 1,112 document--annotator pairs), never inspected or scored before the configuration was frozen, is its held-out confirmation set (protocol committed in \texttt{f2be354}).
\item \textbf{Wikipedia biographies} (real, human gold, held out, second domain): 553 Wikipedia biography summaries annotated with TAB's scheme by the same group \cite{papadopoulou2022bootstrapping} (repository \texttt{anthipapa/textanonymization}, commit \texttt{5df8a0d}; 20 summaries annotated twice, 573 document--annotator pairs). They serve only to test whether the free-text configuration of Section~\ref{sec:prose}, frozen on court judgments, transfers to another genre unchanged; the corpus was added and its manifest frozen before any system was run on it (commit \texttt{365f5ad}).
\item \textbf{CredData} (real credentials, human gold, held out): Samsung's credential dataset \cite{creddata2021} (repository commit \texttt{01e3d13}), lines of real files from 337 open-source repositories that its authors labelled as holding a credential, not holding one, or holding a test value or placeholder. Its builder fetches the pinned commits and replaces every real value with a random value of the same shape; we deleted the raw checkouts as soon as it had done so. One document is one labelled line range; a range with a credential carries each credential value as gold, and a range labelled only as non-credential is gold-free. That gives 64,797 documents, 14,626 gold values and 50,608 gold-free ranges; 1,127 credential labels without a value position were dropped by the loader. Only the candidates proposed by its authors' tools, which included gitleaks and TruffleHog, are labelled, so recall is relative to that candidate set. The protocol and the frozen manifest were committed before any system was run on it (commits \texttt{f160334} and \texttt{8e9e42f}).
\end{itemize}

\begin{table}[t]
\caption{Test splits. Gold-free documents contain no gold span; any detection on them is a false alarm. Clusters are the resampling units of the cluster and two-stage bootstrap (Section~\ref{sec:stats}).}\label{tab:corpora}
\footnotesize
\begin{tabular*}{\textwidth}{@{\extracolsep\fill}lrrrrl}
\toprule
Corpus & Documents & Gold spans & Gold-free & Bytes & Clusters \\
\midrule
presidio-research & 750 & 1,394 & 63 & 63,678 & 200 templates \\
secrets & 313 & 293 & 97 & 46,595 & 96 rules + 5 templates \\
RedactionBench & 100 & 25,592 & 0 & 1,378,267 & 11 domains \\
OpenSSH logs (real) & 2,000 & 2,491 & 148 & 221,218 & 27 templates \\
Linux logs (real) & 2,000 & 1,876 & 407 & 212,487 & 118 templates \\
TAB judgments (real) & 127 & 6,201\footnotemark[1] & 0 & 652,935 & documents \\
TAB held-out (real) & 1,014 & 59,143\footnotemark[1] & 0 & 8,328,627 & documents \\
Wikipedia bios (real) & 553 & 7,527\footnotemark[1] & 0 & 270,206 & documents \\
CredData (real) & 64,797 & 14,626 & 50,608 & 17,943,839 & 337 repositories \\
\bottomrule
\end{tabular*}
\footnotetext[1]{Union over annotators, overlapping mentions merged; TAB is scored against each annotator's gold separately (Section~\ref{sec:metrics}).}
\end{table}

\subsection{Systems and configurations}
\textbf{SafeStream} was run with the Rust core and default tiers (\emph{pattern}), with the opt-in date tier (\emph{+dates}), and with the spaCy \texttt{en\_core\_web\_sm} tier (\emph{+spaCy}). \textbf{Presidio} 2.2.355 \cite{presidio} was run as its pattern recognisers alone (\emph{regex}) and as the full analyser with spaCy \texttt{en\_core\_web\_sm} (\emph{sm}) and \texttt{en\_core\_web\_lg} (\emph{lg}), with spaCy 3.8.13 and models 3.8.0. \textbf{gitleaks} 8.28.0 \cite{gitleaks} and \textbf{TruffleHog} 3.90.8 \cite{trufflehog} were run with default rules and default false-positive filtering; TruffleHog ran with \texttt{--no-verification}, so no example credential was sent to a live service. A sensitivity run repeated both scanners with their false-positive filters disabled (Section~\ref{sec:sens}). Each scanner ran once over a directory with one file per document, and findings were mapped back to character offsets; TruffleHog's composite findings (for example \texttt{user:token}) were credited for each component that appears in the text. Presidio received documents longer than 20,000 characters in line-aligned pieces of at most that length; no gold span in any test split crosses a piece boundary, so the splitting cost Presidio no recall. All SafeStream thresholds are defaults; no threshold was tuned on test data. The RedactionBench dev-selected configurations (Section~\ref{sec:contam}) were chosen on the development half by a rule fixed in advance.

\subsection{Metrics}\label{sec:metrics}
Let $G$ and $R$ be the sets of gold and redacted character positions in a document. Character precision and recall are $|G \cap R| / |R|$ and $|G \cap R| / |G|$, summed over documents, and F1 is their harmonic mean. All metrics are type-agnostic: a value counts as protected whatever label it receives. The \emph{fully redacted rate} is the share of gold spans with every character redacted. The \emph{false-alarm document rate} is the share of gold-free documents with at least one detection. On RedactionBench, \emph{mandatory recall} is character recall over mandatory spans, and the \emph{no-leak document rate} is the share of documents in which every mandatory character is redacted (Section~\ref{sec:noleak} discusses how it relates to recall). TAB is scored as its authors prescribe \cite{pilan2022tab}: recall against each annotator's gold, pooled over document--annotator pairs; precision counting a redacted character as correct if any annotator marked it; and \emph{entity recall} for direct and quasi identifiers, under which an entity counts as protected only if every one of its mentions is fully redacted. TAB's information-weighted precision needs a language model and is not computed.

\subsection{Statistical analysis}\label{sec:stats}
We report 95\% percentile bootstrap intervals with 2,000 resamples and seed 0 \cite{efron1993bootstrap}. All systems evaluated on a corpus are scored on one shared resample matrix, so every difference between two systems is paired; we report the median paired difference and call it significant only when its paired interval excludes zero. Documents are not independent. Secrets positives cluster by rule, presidio-research documents share generator templates, and RedactionBench documents share domains. We therefore use three resampling schemes on the same per-document counts \cite{field2007bootstrapping}: \emph{document} (documents drawn independently), \emph{cluster} (whole clusters drawn with replacement, all their documents kept) and \emph{two-stage} (clusters drawn with replacement, then documents drawn with replacement inside each drawn cluster). Table~\ref{tab:corpora} lists the clusters. Tables~\ref{tab:main} and~\ref{tab:rb} give document-level intervals. Table~S9 of Online Resource~1 repeats the planned comparisons under all three schemes, and we call a difference \emph{robust} when it is significant under all three. RedactionBench has only eleven clusters, so its cluster intervals are themselves imprecise; we read them as a sensitivity analysis rather than as an exact test. No correction for multiple comparisons was applied; the comparisons in the abstract were planned, the ablation comparisons were not. The analysis was not pre-registered.

\subsection{Environment}
Accuracy runs used a 4-vCPU Intel Xeon virtual machine at 2.80\,GHz with 15\,GB RAM, Linux and CPython 3.13. Throughput figures come from separate runs, five per configuration, on the same machine with no other benchmark running; the throughput recorded during accuracy runs is not used. Memory was measured as the peak resident set size of a fresh process per point minus its resident set size after imports, so native allocations are included; it is reported in MiB. Re-running every test-split system for the cluster analysis reproduced each point estimate and document-level interval to three decimals. The held-out results come from commit \texttt{0765dae}, the results after the entropy fix (Section~\ref{sec:fp}) from \texttt{3d2c6ac}, and the pre-exposure RedactionBench results from \texttt{874e105}. On the real-data corpora the evaluated configurations ran at \texttt{3d2c6ac} with the pure-Python backend, whose output equals the Rust core's; the user-name tier and the current default profile ran at the current code, and the free-text configuration at \texttt{9c10975}, where it was frozen. All commits are on the main branch of the public repository.
